% --- Archon physics formalization source begin ---
% archon:physics
% archon:covers PhyXMiniProblems/problem_phyx_mini_0013.lean
% archon:source-report reports/phyx_mini/problem_phyx_mini_0013.source.json

\chapter{Physics problem phyx\_mini\_0013}
\label{ch:PhyXMiniProblems_problem_phyx_mini_0013}

\paragraph{Problem source.}
\#\# Physical scenario

A digital video disc (DVD) records information in a spiral track approximately \textbackslash{}( 1\textbackslash{} \textbackslash{}mu\textbackslash{}text\{m\} \textbackslash{}) wide. The track consists of a series of pits in the information layer that scatter light from a laser beam sharply focused on them. The laser shines in from below through transparent plastic of thickness \textbackslash{}( t = 1.20\textbackslash{} \textbackslash{}text\{mm\} \textbackslash{}) and index of refraction 1.55 as shown in figure. Assume the width of the laser beam at the information layer must be \textbackslash{}( a = 1.00\textbackslash{} \textbackslash{}mu\textbackslash{}text\{m\} \textbackslash{}) to read from only one track and not from its neighbors. Assume the width of the beam as it enters the transparent plastic is \textbackslash{}( w = 0.700\textbackslash{} \textbackslash{}text\{mm\} \textbackslash{}). A lens makes the beam converge into a cone with an apex angle \textbackslash{}( 2\textbackslash{}theta\_1 \textbackslash{}) before it enters the DVD.

\#\# Question

Find the incidence angle \textbackslash{}( \textbackslash{}theta\_1 \textbackslash{}) of the light at the edge of the conical beam. This design is relatively immune to small dust particles degrading the video quality.

\#\# Answer choices

- A: A: \textbackslash{}( 20.6\textasciicircum{}\textbackslash{}circ \textbackslash{})

- B: B: \textbackslash{}( 18.3\textasciicircum{}\textbackslash{}circ \textbackslash{})

- C: C: \textbackslash{}( 25.7\textasciicircum{}\textbackslash{}circ \textbackslash{})

- D: D: \textbackslash{}( 21.3\textasciicircum{}\textbackslash{}circ \textbackslash{})

\#\# Figure caption (auxiliary; use the image as primary evidence)

The image is a diagram illustrating light refraction through a plastic layer with the refractive index, \textbackslash{}(n = 1.55\textbackslash{}).

\#\#\# Key Elements:

1. **Layers:**
   - **Information Layer:** Positioned at the top with a horizontal width.
   - **Plastic Layer**: Below the information layer, labeled with \textbackslash{}(n = 1.55\textbackslash{}) and separated by dashed lines on the sides.

2. **Dimensions:**
   - **\textbackslash{}(b\textbackslash{}):** Half the width of the information layer on either side.
   - **\textbackslash{}(a\textbackslash{}):** Central width of the information layer.
   - **\textbackslash{}(t\textbackslash{}):** Total vertical height of the plastic layer.
   - **\textbackslash{}(w\textbackslash{}):** Base width between the exiting light rays in the air.

3. **Light Rays:**
   - Two light rays are shown entering from the bottom in the air, refracting inside the plastic layer, and exiting towards the top.
   - **Angles:** 
     - \textbackslash{}(\textbackslash{}theta\_1\textbackslash{}): Angle of incidence in air.
     - \textbackslash{}(\textbackslash{}theta\_2\textbackslash{}) and \textbackslash{}(\textbackslash{}theta'\_2\textbackslash{}): Angles of refraction inside the plastic.

4. **Material:**
   - **Air:** Denoted below the plastic layer where the initial light rays originate.

\#\#\# Relationships:
The image illustrates the refraction process: light enters the plastic layer at \textbackslash{}(\textbackslash{}theta\_1\textbackslash{}), changes direction due to the refractive index, travels through the plastic at \textbackslash{}(\textbackslash{}theta\_2\textbackslash{}) and \textbackslash{}(\textbackslash{}theta'\_2\textbackslash{}), and exits back into air. The geometry and refractive properties are annotated to demonstrate these relationships.

\#\# Dataset metadata

Geometrical Optics; Physical Model Grounding Reasoning, Spatial Relation Reasoning

\paragraph{Recorded answer/context.}
C: C: \textbackslash{}( 25.7\textasciicircum{}\textbackslash{}circ \textbackslash{})

\paragraph{Figure/image path.}
/root/proposal\_for\_physic/science-mango/phyx\_mini\_run/phyx\_data/test\_image/13.png

\paragraph{Formalization target.}
create a compiling Lean file with sorry bodies at `PhyXMiniProblems/problem\_phyx\_mini\_0013.lean`. The Lean declarations must preserve the physical quantities, dimensions or dimensional roles, figure labels, governing-law hypotheses, and final relation expressed by this problem.
Use Mathlib/Physlib names found through LeanExplore where available. If a physics API is missing, introduce faithful local abstractions rather than scalar placeholder aliases.

\begin{theorem}[Physics formalization target]
\label{thm:physics:phyx_mini_0013:target}
\lean{PhyXMiniProblems.ProblemPhyXMini0013.incidenceAngle_is_answer_C}
\uses{def:physics:phyx-mini-0013:phyxminiproblems-problemphyxmini0013-dvdlasersetup, def:physics:phyx-mini-0013:phyxminiproblems-problemphyxmini0013-dvdlaserfigurereadouts, def:physics:phyx-mini-0013:phyxminiproblems-problemphyxmini0013-dvdlasergoverninglaws, def:physics:phyx-mini-0013:phyxminiproblems-problemphyxmini0013-matchesanswertonearesttenthdegree}
For the DVD geometry with plastic thickness \(1.20\,\mathrm{mm}\), track
width \(1.00\,\mu\mathrm m\), entry width \(0.700\,\mathrm{mm}\), and
dimensionless plastic index \(1.55\), the edge-ray incidence angle in air
rounds to \(25.7^\circ\), answer C.
\end{theorem}
\begin{proof}
Convert all lateral readouts to one length unit.  The two right-triangle
relations determine the acute plastic edge angles from the half-width changes
over the slab thickness.  Substitution in the two Snell equations then
determines the common acute incidence angle; certified trigonometric bounds
place it within \(0.05^\circ\) of \(25.7^\circ\).
\end{proof}

% --- Archon declaration topology begin ---
\section{Declaration topology}
\begin{definition}[Length Quantity]
  \label{def:physics:phyx-mini-0013:phyxminiproblems-problemphyxmini0013-lengthquantity}
  \lean{PhyXMiniProblems.ProblemPhyXMini0013.LengthQuantity}
  A physical length, represented independently of the unit system used to read it
\end{definition}
\begin{proof}
  This is the declared model definition of Length Quantity.
\end{proof}
\begin{definition}[value In Meters]
  \label{def:physics:phyx-mini-0013:phyxminiproblems-problemphyxmini0013-valueinmeters}
  \lean{PhyXMiniProblems.ProblemPhyXMini0013.valueInMeters}
  \uses{def:physics:phyx-mini-0013:phyxminiproblems-problemphyxmini0013-lengthquantity}
  The scalar value of a physical length when the chosen unit system is SI meters
\end{definition}
\begin{proof}
  This is the declared model definition of value In Meters.
\end{proof}
\begin{definition}[Optical Medium]
  \label{def:physics:phyx-mini-0013:phyxminiproblems-problemphyxmini0013-opticalmedium}
  \lean{PhyXMiniProblems.ProblemPhyXMini0013.OpticalMedium}
  The two optical media separated by the lower face of the DVD
\end{definition}
\begin{proof}
  This is the declared model definition of Optical Medium.
\end{proof}
\begin{definition}[DVDLaser Setup]
  \label{def:physics:phyx-mini-0013:phyxminiproblems-problemphyxmini0013-dvdlasersetup}
  \lean{PhyXMiniProblems.ProblemPhyXMini0013.DVDLaserSetup}
  \uses{def:physics:phyx-mini-0013:phyxminiproblems-problemphyxmini0013-lengthquantity, def:physics:phyx-mini-0013:phyxminiproblems-problemphyxmini0013-opticalmedium}
  The physical quantities and figure labels in the symmetric DVD laser-beam cross-section. The angles are real-valued radian readouts measured from the interface normal. The two in-plastic angles retain the figure's left and right labels theta\_2 and theta'\_2
\end{definition}
\begin{proof}
  This is the declared model definition of DVDLaser Setup.
\end{proof}
\begin{definition}[Snell Law At Interface]
  \label{def:physics:phyx-mini-0013:phyxminiproblems-problemphyxmini0013-snelllawatinterface}
  \lean{PhyXMiniProblems.ProblemPhyXMini0013.SnellLawAtInterface}
  Snell's law at a planar interface. Refractive indices are dimensionless and angles are radian magnitudes from the local normal
\end{definition}
\begin{proof}
  This is the declared model definition of Snell Law At Interface.
\end{proof}
\begin{definition}[DVDLaser Figure Readouts]
  \label{def:physics:phyx-mini-0013:phyxminiproblems-problemphyxmini0013-dvdlaserfigurereadouts}
  \lean{PhyXMiniProblems.ProblemPhyXMini0013.DVDLaserFigureReadouts}
  \uses{def:physics:phyx-mini-0013:phyxminiproblems-problemphyxmini0013-valueinmeters, def:physics:phyx-mini-0013:phyxminiproblems-problemphyxmini0013-dvdlasersetup}
  Problem-statement and figure readouts. Lengths are read in SI meters, so the displayed micrometer and millimeter values are converted exactly to powers of ten. The width partition and cone-angle relations are geometric annotations from the symmetric figure; neither states the requested numerical incidence angle
\end{definition}
\begin{proof}
  This is the declared model definition of DVDLaser Figure Readouts.
\end{proof}
\begin{definition}[DVDLaser Governing Laws]
  \label{def:physics:phyx-mini-0013:phyxminiproblems-problemphyxmini0013-dvdlasergoverninglaws}
  \lean{PhyXMiniProblems.ProblemPhyXMini0013.DVDLaserGoverningLaws}
  \uses{def:physics:phyx-mini-0013:phyxminiproblems-problemphyxmini0013-opticalmedium, def:physics:phyx-mini-0013:phyxminiproblems-problemphyxmini0013-dvdlasersetup, def:physics:phyx-mini-0013:phyxminiproblems-problemphyxmini0013-snelllawatinterface}
  The governing geometrical-optics relations for the two edge rays. The right-triangle equations express b = t tan(theta\_2) in any chosen unit system, and the two Snell-law fields describe refraction from air into the plastic. Acute-angle hypotheses select the physical branches of the trigonometric equations
\end{definition}
\begin{proof}
  This is the declared model definition of DVDLaser Governing Laws.
\end{proof}
\begin{definition}[Answer Choice]
  \label{def:physics:phyx-mini-0013:phyxminiproblems-problemphyxmini0013-answerchoice}
  \lean{PhyXMiniProblems.ProblemPhyXMini0013.AnswerChoice}
  The four incidence-angle choices printed with the problem
\end{definition}
\begin{proof}
  This is the declared model definition of Answer Choice.
\end{proof}
\begin{definition}[answer Angle Degrees]
  \label{def:physics:phyx-mini-0013:phyxminiproblems-problemphyxmini0013-answerangledegrees}
  \lean{PhyXMiniProblems.ProblemPhyXMini0013.answerAngleDegrees}
  \uses{def:physics:phyx-mini-0013:phyxminiproblems-problemphyxmini0013-answerchoice}
  The degree readout printed beside each answer choice
\end{definition}
\begin{proof}
  This is the declared model definition of answer Angle Degrees.
\end{proof}
\begin{definition}[degrees To Radians]
  \label{def:physics:phyx-mini-0013:phyxminiproblems-problemphyxmini0013-degreestoradians}
  \lean{PhyXMiniProblems.ProblemPhyXMini0013.degreesToRadians}
  Convert a real-valued degree readout to radians
\end{definition}
\begin{proof}
  This is the declared model definition of degrees To Radians.
\end{proof}
\begin{definition}[Matches Answer To Nearest Tenth Degree]
  \label{def:physics:phyx-mini-0013:phyxminiproblems-problemphyxmini0013-matchesanswertonearesttenthdegree}
  \lean{PhyXMiniProblems.ProblemPhyXMini0013.MatchesAnswerToNearestTenthDegree}
  \uses{def:physics:phyx-mini-0013:phyxminiproblems-problemphyxmini0013-answerchoice, def:physics:phyx-mini-0013:phyxminiproblems-problemphyxmini0013-answerangledegrees, def:physics:phyx-mini-0013:phyxminiproblems-problemphyxmini0013-degreestoradians}
  A radian angle rounds to a displayed one-decimal-place answer when it lies within half of 0.1 degrees of that choice's degree readout
\end{definition}
\begin{proof}
  This is the declared model definition of Matches Answer To Nearest Tenth Degree.
\end{proof}
% --- Archon declaration topology end ---
% --- Archon physics formalization source end ---
